\section{Introdução}
\label{sec:introducao}

O monitoramento ambiental por \ac{RSSF} permite a medição contínua e distribuída de variáveis relacionadas ao microclima, à qualidade do ar, às condições do solo, à qualidade da água e à ocorrência de eventos naturais \cite{HART2006177}.
Para viabilizar essa observação, uma aplicação de \ac{RSSF} para monitoramento ambiental consiste em uma rede distribuída de nós sensores que realizam medições \textit{in situ} e encaminham os dados, por enlaces sem fio, a um nó sorvedouro ou gateway, que os disponibiliza a um servidor de armazenamento e processamento \cite{5597912,AKYILDIZ2002393}.

Essa configuração reduz a necessidade de cabeamento de comunicação, facilita a instalação de múltiplos pontos de medição e possibilita observações de longa duração em áreas extensas, remotas ou de difícil acesso \cite{HART2006177,5597912}.
Como resultado, uma das principais vantagens das \acp{RSSF} no monitoramento ambiental é combinar elevada resolução espacial, mediante medições em múltiplos locais, com elevada resolução temporal, obtida por amostragens repetidas e coordenadas \cite{HART2006177,5597912}.

Essa capacidade foi demonstrada em implantações clássicas em habitats naturais, nas quais uma \ac{RSSF} registrou temperatura, umidade, pressão atmosférica e luminosidade, reduzindo a necessidade de visitas frequentes e a consequente perturbação do local \cite{10.1145/570738.570751}.
Para que os registros resultantes tenham utilidade científica, devem ser consideradas a qualidade das medições, a consistência das marcas temporais, a localização dos nós, a continuidade das séries e a confiabilidade da transmissão \cite{5597912,YICK20082292,10.1145/1031495.1031521}.